\chapter{Cyclic traces and wheel corrections}
\label{chap:vol2-cyclic-traces-wheels}

A matrix trace remembers the size of the matrix:
$\operatorname{Tr}(I_N)=N$.
Invariant quantization also remembers a scalar. For the quadratic
$\mathfrak{sl}_2$ invariant, the Duflo operator adds
$\hbar^2/2$ to the symmetrized Casimir.
These scalars matter when matrices are stabilized or Lie algebras
are completed. Removing them can destroy compatibility with
brackets or continuity.

We first construct trace maps on cyclic chains and specify the
rank transitions they preserve. We then compute the finite wheel
corrections and examine two different infinite-dimensional
coefficient spaces. The cyclic orientation of
Chapter~\ref{ch:vol2-derived-hamiltonian-actions} remains a
chain-level structure throughout. A scalar state requires its own
functional.

\section{Matrix trace on the mixed complex}
\label{sec:vol2-matrix-mixed-trace}

Let $B$ be a unital associative complex algebra and put
$\overline B=B/\mathbb C1$.
Its normalized Hochschild chains are
\[
 C_n(B)=B\otimes\overline B^{\otimes n},
 \qquad a_0[a_1|\cdots|a_n]\in C_n(B).
\]
Place $C_n(B)$ in cohomological degree $-n$.
For $n>0$ the Hochschild boundary is
\begin{align*}
 b(a_0[a_1|\cdots|a_n])
 ={}&a_0a_1[a_2|\cdots|a_n]\\
 &+\sum_{j=1}^{n-1}(-1)^j
      a_0[a_1|\cdots|a_ja_{j+1}|\cdots|a_n]\\
 &+(-1)^na_na_0[a_1|\cdots|a_{n-1}],
\end{align*}
and $b=0$ on $C_0(B)$.
The Connes operator is
\[
 \mathsf B(a_0[a_1|\cdots|a_n])
 =\sum_{j=0}^n(-1)^{nj}
       1[a_j|\cdots|a_n|a_0|\cdots|a_{j-1}].
\]
Units in bar slots are zero. The identities
$b^2=\mathsf B^2=b\mathsf B+\mathsf Bb=0$ give the mixed
complex. With $|\zeta|=2$, its negative-cyclic complex is
\[
 CC^-(B)=\bigl(C_\bullet(B)[[\zeta]],b+\zeta\mathsf B\bigr).
\]
Here the completion is in $\zeta$. In particular, it does not
complete the coefficients of $B$.

For a second unital algebra $D$, define
\begin{equation}\label{eq:vol2-matrix-chain-trace}
 \mathsf T_N\bigl(M_0[M_1|\cdots|M_n]\bigr)
 =
 \sum_{i_0,\ldots,i_n=1}^N
 (M_0)_{i_0i_1}
 [(M_1)_{i_1i_2}|\cdots|(M_n)_{i_ni_0}],
\end{equation}
from matrix chains to chains of $D$.
For $n=0$ the formula is the ordinary matrix trace.
The indices form a closed path, so a change of the first index
does not change the sum.

\begin{proposition}\label{prop:vol2-mixed-matrix-trace}
The map $\mathsf T_N$ is a morphism of mixed complexes
$C_\bullet(\operatorname{Mat}_N D)\to C_\bullet(D)$.
It therefore extends continuously to negative-cyclic complexes.
It is natural in $D$ and additive under block sums of
representations. Its value on the unit in degree zero is $N1$.
\end{proposition}
\begin{proof}
A scalar identity in a bar slot forces the two adjacent
indices to agree. The corresponding output has a scalar unit
in that bar slot and is zero. Thus the formula is defined on
normalized chains.

An internal Hochschild face multiplies two adjacent matrices.
Expanding the product sums precisely over the shared index
in \eqref{eq:vol2-matrix-chain-trace}.
The first face has the same interpretation at the outer
coefficient. For the last face, reindex the closed path
cyclically. These identities prove
$b\mathsf T_N=\mathsf T_Nb$, with the displayed face signs.

In a Connes term the leading identity matrix forces its
adjacent indices to coincide. After deleting that repeated
index, the remaining sum is the cyclic rotation of the same
closed path. The rotation has the same sign $(-1)^{nj}$ on
both sides. Hence
$\mathsf B\mathsf T_N=\mathsf T_N\mathsf B$.
Coefficientwise application now commutes with
$b+\zeta\mathsf B$ and with reduction modulo every power of
$\zeta$.

For block-diagonal matrices, a nonzero closed path lies
entirely in one block. The trace of a block sum is consequently
the sum of the two chain traces. Naturality follows by
applying an algebra map to every matrix entry. Finally the
ordinary trace of the identity is $N$.
\end{proof}

For dg algebras use the raw totalization and graded cyclic
signs of
\eqref{eq:vol2-raw-hochschild-differential} and
\eqref{eq:vol2-raw-connes-operator}.
The same formula is a mixed-complex map: matrix indices
have degree zero, the internal differential acts entrywise,
and every cyclic permutation has the same Koszul sign
before and after the index sum.

This gives a concrete image of the orientation already
constructed on the free resolution.
Let
\[
 Q=\mathbb C\langle q,p,t\rangle,\qquad
 |t|=-1,\qquad dt=qp-pq,
\]
and let $R_N$ be its derived representation algebra:
the graded-commutative algebra on the entries of $X,Y,\Psi$,
with
$dX=dY=0$ and $d\Psi=XY-YX$.
Universal substitution is a dg algebra map
$Q\to\operatorname{Mat}_N(R_N)$.
The trace of its negative-cyclic orientation is
\begin{equation}\label{eq:vol2-traced-orientation}
 \omega_N=
 \sum_{i,j}
 \bigl(1[X_{ij}|Y_{ji}]-1[Y_{ij}|X_{ji}]\bigr)
 -\sum_i1[\Psi_{ii}].
\end{equation}
It is closed for both mixed differentials, by
Proposition~\ref{prop:vol2-mixed-matrix-trace}.
The relation term has total degree $-2$, just as the
two-bar terms do.
The matrix trace also commutes with entrywise derivations,
so it transports the coherent classical action and
orientation of
Theorem~\ref{thm:vol2-coherent-cyclic-orientation}.
These are constructions on derived representation
coefficients. The full quantum BRST and supported
field-theoretic comparisons retain their additional maps.

\section{Rank and scalar reduction}
\label{sec:vol2-reduced-stable-trace}

The degree-zero cyclic space of $B$ is
\[
 \mathfrak c_B=HC_0(B)=B/[B,B],
\]
where $[B,B]$ is the linear span of commutators.
Let $\operatorname{Rep}_NB$ represent unital algebra maps
$B\to\operatorname{Mat}_N(D)$ for commutative complex
algebras $D$, and put
\[
 \mathcal R_N=\mathcal O(\operatorname{Rep}_NB)^{GL_N}.
\]
The ordinary trace is the linear map
\[
 \operatorname{Tr}_N:\mathfrak c_B\longrightarrow\mathcal R_N,
 \qquad
 \operatorname{Tr}_N([b])(\rho)=\operatorname{Tr}(\rho(b)).
\]
The trace of a commutator is zero and conjugation preserves
the trace, proving that this map has the stated source
and target. It is generally not multiplicative.
Even on scalars, its value at $1$ is $N$.

Choose an augmentation $\epsilon:B\to\mathbb C$.
It specifies a rank transition
\[
 s_N:\operatorname{Rep}_NB\longrightarrow\operatorname{Rep}_{N+1}B,
 \qquad \rho\longmapsto\rho\oplus\epsilon.
\]
Define
\[
 \overline{\mathfrak c}_B=B/(\mathbb C1+[B,B]),
 \qquad
 t_N(\overline b)(\rho)=\operatorname{Tr}(\rho(b))-N\epsilon(b).
\]
Let $\mathfrak m_N$ be the ideal in $\mathcal R_N$ of functions
vanishing at $b\mapsto\epsilon(b)I_N$, and set
\[
 \widehat{\mathcal R}_N
   =\varprojlim_{r\geq1}\mathcal R_N/\mathfrak m_N^r.
\]
If $J$ is the ideal generated by
$\overline{\mathfrak c}_B$ in
$\operatorname{Sym}_{\mathbb C}\overline{\mathfrak c}_B$,
write
\[
 \widehat{\operatorname{Sym}}_J\overline{\mathfrak c}_B
 =\varprojlim_{r\geq1}
   \operatorname{Sym}_{\mathbb C}\overline{\mathfrak c}_B/J^r.
\]

\begin{proposition}\label{prop:vol2-reduced-trace-completion}
The maps $t_N$ are well defined, have image in
$\mathfrak m_N$, and satisfy $s_N^*t_{N+1}=t_N$.
They extend to a continuous homomorphism of commutative algebras
\[
 \widehat{\operatorname{Sym}}_J\overline{\mathfrak c}_B
 \longrightarrow \varprojlim_N\widehat{\mathcal R}_N.
\]
All maps are natural for augmentation-preserving algebra maps.
\end{proposition}
\begin{proof}
Both terms defining $t_N$ vanish on commutators, and their
values on a scalar multiple of $1$ cancel.
Their values also cancel at the augmentation representation.
Moreover,
\[
 s_N^*t_{N+1}(\overline b)
 =\operatorname{Tr}_N([b])+\epsilon(b)-(N+1)\epsilon(b)
 =t_N(\overline b).
\]
Since $s_N$ takes the augmentation representation to the
augmentation representation,
$s_N^*(\mathfrak m_{N+1}^r)\subseteq\mathfrak m_N^r$.
The symmetric-algebra extension of $t_N$ sends $J^r$ into
$\mathfrak m_N^r$. It therefore defines a map on each pair of
quotients and hence on their inverse limits.
The stabilization identity makes these completed maps
compatible in $N$. These formulas also prove naturality.
\end{proof}

The subtraction exists on cyclic chains as well.
For a representation $\rho:B\to\operatorname{Mat}_N(D)$,
subtract $N$ times the map induced by
$B\xrightarrow{\epsilon}\mathbb C\to D$ from
$\mathsf T_N\rho_*$.
Both summands are mixed-complex maps.
Their difference vanishes on the scalar subcomplex and is
unchanged by $\rho\mapsto\rho\oplus\epsilon$.
This chain-level reduction has a chosen augmentation as
part of its data.

A Poisson comparison requires more.
Suppose $\mathfrak c_B$ is supplied with a Lie bracket and
$\mathcal R_N$ with Poisson brackets satisfying
\begin{equation}\label{eq:vol2-trace-bracket-datum}
 \{\operatorname{Tr}_N\beta,\operatorname{Tr}_N\gamma\}_N
 =\operatorname{Tr}_N[\beta,\gamma]_{\mathrm{cyc}}.
\end{equation}
Then the ordinary trace is a Lie map.
A Hamiltonian action through it also requires a Lie map
$j:\mathfrak h\to\mathfrak c_B$.
Neither datum is attached to an arbitrary associative
algebra by its definition.
For the free algebra on two generators, the necklace bracket
and trace moment map provide
\eqref{eq:vol2-trace-bracket-datum};
the polynomial-to-necklace comparison still has the cubic
defect discussed in Chapter~\ref{ch:vol2-derived-hamiltonian-actions}.

Scalar reduction need not preserve this bracket.
For $B=\mathbb C[x,y]$ with $\{x,y\}=1$ and augmentation
at the origin, evaluation on $N$ points gives
\[
 \{\operatorname{Tr}_Nx,\operatorname{Tr}_Ny\}=N,
 \qquad t_N(\overline1)=0.
\]
Thus Proposition~\ref{prop:vol2-reduced-trace-completion}
asserts a commutative algebra map.
A reduced Poisson map needs a bracket-compatible reduction,
for example a selected Hamiltonian sector on which the
augmentation of every bracket vanishes.

\section{The stable multisymmetric algebra}
\label{sec:vol2-stable-multisymmetric}

The rank dependence can be retained explicitly.
For $N$ points of the affine symplectic plane put
\[
 S_N=\mathbb C[x_1,y_1,\ldots,x_N,y_N]^{\mathfrak S_N},
 \qquad
 p_{ab}^{(N)}=\sum_{i=1}^Nx_i^ay_i^b.
\]
Give $x_i,y_i$ degree one and, for $a+b>0$, give the
formal generator $p_{ab}$ degree $a+b$.
Evaluation at an extra point $(0,0)$ defines
$S_{N+1}\to S_N$.

\begin{theorem}\label{thm:vol2-stable-multisymmetric-ring}
In each fixed degree $d$, evaluation gives an isomorphism
\[
 \mathbb C[p_{ab}\mid a,b\geq0,\ a+b>0]_d
   \xrightarrow{\ \sim\ }(S_N)_d
 \qquad (N\geq d).
\]
Writing $P$ for this graded polynomial algebra, one has
\[
 \widehat P_{\mathrm{wt}}=\prod_{d\geq0}P_d
 \cong
 \varprojlim_N\prod_{d\geq0}(S_N)_d
\]
as complete commutative algebras.
\end{theorem}
\begin{proof}
A monomial orbit is specified by a finite multiset
$\lambda$ of nonzero pairs $(a,b)$, one pair for each
occupied point. Let $m_\lambda^{(N)}$ be the sum of the
distinct monomials in this orbit.
These sums form a basis when $\lambda$ has at most $N$
members. In total degree $d$, it has at most $d$ members.
The basis is therefore unchanged for $N\geq d$, and
evaluation at the added origin preserves every basis element.

For $\lambda=\{\alpha_1,\ldots,\alpha_r\}$ expand
$p_{\alpha_1}^{(N)}\cdots p_{\alpha_r}^{(N)}$.
Terms using distinct point indices give
\[
 \left(\prod_\alpha m_\alpha!\right)m_\lambda^{(N)},
\]
where $m_\alpha$ is the multiplicity of $\alpha$ in
$\lambda$. Every other term uses fewer than $r$ points,
merging the pairs assigned to the same point.
Ordering orbit sums by the number of occupied points
gives a triangular change of basis with nonzero diagonal.
It proves both generation and independence in degree $d$.

Multiplication adds total degrees. Completing degreewise
therefore gives $\widehat P_{\mathrm{wt}}$.
A compatible family in the right-hand limit determines,
in each degree, the unique stable coefficient just
constructed. Conversely those coefficients determine the
family. These inverse maps preserve multiplication and
the inverse-limit topologies.
\end{proof}

The completion in this theorem is specified by total degree.
It differs from the $J$-adic completion in which every
$p_{ab}$ has filtration one.
For example,
$\sum_{a\geq1}p_{a0}$ belongs to
$\widehat P_{\mathrm{wt}}$, but does not define an element
of the $J$-adic completion: its reduction modulo $J^2$
would be an infinite linear combination in a vector space
whose elements are finite linear combinations.
Nor does the theorem replace the target by the inverse
limit of the uncompleted rings $S_N$.
The series $\sum_{r\geq0}p_{10}^r$ already has a
nonpolynomial image at $N=1$.

These degree-zero trace functions do not record
the relation and gauge cochain degrees of a
derived matrix quotient. A comparison with a
derived centre requires representation and trace
maps on the specified complexes.

For the Poisson bracket, retain a central generator
$r=p_{00}$ and put
\[
 P^{\mathrm{rk}}=\mathbb C[r,p_{ab}\mid a+b>0].
\]
Define on generators
\begin{equation}\label{eq:vol2-rank-poisson}
 \{p_{ab},p_{cd}\}=(ad-bc)p_{a+c-1,b+d-1},
 \qquad \{r,-\}=0.
\end{equation}
A term with a negative index is zero; in that case its
coefficient is already zero.

\begin{proposition}\label{prop:vol2-rank-poisson}
Equation~\eqref{eq:vol2-rank-poisson} defines a Poisson
bracket. The evaluations
$P^{\mathrm{rk}}\to S_N$, with $r\mapsto N$, are Poisson.
Block sum is described by the Poisson coproduct
\[
 \Delta p_{ab}=p_{ab}\otimes1+1\otimes p_{ab},
 \qquad
 \Delta r=r\otimes1+1\otimes r.
\]
Origin stabilization is not Poisson on the full $S_N$.
On the subalgebras generated by $p_{ab}^{(N)}$ with
$a+b\geq2$, it is Poisson, and their stable algebra is
$\mathbb C[p_{ab}\mid a+b\geq2]$ with the same bracket.
\end{proposition}
\begin{proof}
The polynomial Poisson algebra $\mathbb C[x,y]$, regarded
only as a Lie algebra, has basis $x^ay^b$ including $1$.
Its symmetric algebra is $P^{\mathrm{rk}}$, where the
Lie element $1$ becomes the generator $r$, not the
multiplicative unit. The Lie--Poisson construction extends
its bracket by Leibniz and proves Jacobi.

Direct differentiation in $S_N$ gives
\[
 \{p_{ab}^{(N)},p_{cd}^{(N)}\}
 =(ad-bc)p_{a+c-1,b+d-1}^{(N)},\qquad
 p_{00}^{(N)}=N.
\]
This proves the evaluation assertion.
The primitive coproduct preserves a Lie bracket on
generators, so its extension is Poisson.
It also equals the power-sum formula for a block sum.

In rank $N+1$ the bracket of $p_{10}$ and $p_{01}$ is
$N+1$. Its pullback remains $N+1$, whereas the bracket
of the pullbacks is $N$. This proves the failure.
If both input degrees are at least two, the output
degree is at least two. The central rank never occurs,
and setting the last point to zero preserves every
displayed bracket. The triangular basis argument of
Theorem~\ref{thm:vol2-stable-multisymmetric-ring}, restricted
to these generators, proves the final assertion.
\end{proof}

\section{Stable matrix Lie homology}
\label{sec:vol2-stable-lie-homology}

There is a second stabilization, with a different direction.
Let
\[
 \mathfrak{gl}_\infty B=\varinjlim_N\mathfrak{gl}_N B
\]
under block-zero Lie embeddings, and use trivial
coefficients $\mathbb C$ for its Chevalley homology.
This construction requires no augmentation of $B$.
For clarity, ordinary cyclic homology is the homology
of the cyclic complex
\[
 CC(B)=
 \bigl(C_\bullet(B)((\zeta))/\zeta C_\bullet(B)[[\zeta]],
       b+\zeta\mathsf B\bigr),
 \qquad HC_n(B)=H^{-n}CC(B).
\]

The Loday--Quillen--Tsygan theorem identifies its
primitive part:
\[
 \operatorname{Prim}H_n^{\mathrm{CE}}(\mathfrak{gl}_\infty B;\mathbb C)
 \cong HC_{n-1}(B).
\]
Equivalently, if $HC_{n-1}(B)$ is placed in homological
degree $n$, it gives an isomorphism of graded Hopf algebras
\begin{equation}\label{eq:vol2-lqt-hopf}
 H_*^{\mathrm{CE}}(\mathfrak{gl}_\infty B;\mathbb C)
 \cong
 \operatorname{Sym}_{\mathrm{gr}}
       \left(\bigoplus_{n\geq1}HC_{n-1}(B)\right).
\end{equation}
Here $\operatorname{Sym}_{\mathrm{gr}}$ uses the Koszul
sign rule. This is the characteristic-zero theorem for
unital associative algebras of Loday and Quillen
\cite{LodayQuillen1984}*{Theorem~6.2 and Proposition~6.6}.
The product comes from block sum and the coproduct
from the diagonal Lie map.

The matrix-chain trace in
\eqref{eq:vol2-matrix-chain-trace} is also the trace
appearing in this theorem. The invariant-theory
description separates a permutation into its cycles:
one cycle gives a primitive cyclic chain, and disjoint
cycles give its graded symmetric products.
The theorem concerns this stable Lie homology, not the
inverse system of representation-coordinate rings.

For $B=\mathbb C$, normalized Hochschild chains are
concentrated in degree zero. Thus $HC_{2j}(\mathbb C)=\mathbb C$
and $HC_{2j+1}(\mathbb C)=0$, and
\[
 H_*^{\mathrm{CE}}(\mathfrak{gl}_\infty\mathbb C;\mathbb C)
 \cong\Lambda(u_1,u_3,u_5,\ldots),\qquad |u_{2j+1}|=2j+1.
\]
These classes remain although the scalar-reduced cyclic
space in Proposition~\ref{prop:vol2-reduced-trace-completion}
is zero.

\section{Cyclic functionals on the Weyl algebra}
\label{sec:vol2-weyl-cyclic-functional}

Quantization further separates cyclic chains from a scalar
trace. Let $R=\mathbb C[\hbar]$ and let $W_\hbar$ be the
polynomial Weyl algebra generated by $x,y$ with
$[x,y]=\hbar$.
Denote by $[W_\hbar,W_\hbar]_R$ the $R$-linear span
of commutators.

\begin{proposition}\label{prop:vol2-weyl-cyclic-quotient}
One has
\[
 [W_\hbar,W_\hbar]_R=\hbar W_\hbar,\qquad
 HH_0^R(W_\hbar)=W_\hbar/\hbar W_\hbar.
\]
Every $R$-linear cyclic functional $W_\hbar\to R$ is zero.
Every continuous $\mathbb C[[\hbar]]$-linear cyclic
functional on its $\hbar$-adic completion with values
in $\mathbb C[[\hbar]]$ is also zero.
\end{proposition}
\begin{proof}
Reduction modulo $\hbar$ is commutative, so every
commutator belongs to $\hbar W_\hbar$.
For every normal monomial,
\[
 [x,x^ay^{b+1}]=\hbar(b+1)x^ay^b.
\]
Since $b+1$ is invertible in $\mathbb C$, these
commutators span $\hbar W_\hbar$.
The Hochschild boundary in degree one is a commutator,
which gives the quotient formula.

If $\tau$ is cyclic and $R$-linear, then
$\hbar\tau(x^ay^b)=0$ for every $a,b$.
The target has no $\hbar$-torsion, so all these values
vanish. Continuity extends the same argument from
the dense polynomial algebra to its completion.
\end{proof}

The quotient in the proposition is nonzero and has
$\hbar$ acting as zero. It is the universal cyclic
target in degree zero, rather than a torsion-free
scalar ring. The negative-cyclic orientation of the
resolved plane has a different degree and universal
property. Neither it nor a prepared boundary state
is removed by the proposition.

\section{Finite formality and the Duflo map}
\label{sec:vol2-finite-duflo}

Let $\mathfrak g$ be a finite-dimensional complex Lie
algebra. Its linear Poisson tensor on $\mathfrak g^*$ is
$\pi_{\mathfrak g}$.
The polynomial polyvector complex, in unshifted
cochain degrees, is
\[
 \mathcal P_{\mathfrak g}
 =
 \left(
   \bigl(\operatorname{Sym}\mathfrak g
          \otimes\Lambda^\bullet\mathfrak g^*\bigr)[[\hbar]],
   \hbar d_{\mathrm{CE}}
 \right).
\]
The differential is
$[\hbar\pi_{\mathfrak g},-]_{\mathrm{Sch}}$.
Let $A_\hbar$ be the $\hbar$-adic completion of the
Rees enveloping algebra with relations
$xy-yx=\hbar[x,y]_{\mathfrak g}$.
Coefficients remain polynomial at each parameter order.

Fix the affine formality morphism $\mathcal U$ with
first component the normalized
Hochschild--Kostant--Rosenberg map:
\[
 X_1\wedge\cdots\wedge X_p
 \longmapsto
 \frac1{p!}\sum_{\sigma\in\mathfrak S_p}
       \operatorname{sgn}(\sigma)
       X_{\sigma(1)}\otimes\cdots\otimes X_{\sigma(p)}.
\]
We use Kontsevich's affine theorem
\cite{KontsevichDQ}*{Section~6.4}.
Twisting by $\hbar\pi_{\mathfrak g}$ gives a product
$\star$ on $\operatorname{Sym}\mathfrak g[[\hbar]]$.
There is an algebra isomorphism
\[
 T:(\operatorname{Sym}\mathfrak g[[\hbar]],\star)
       \longrightarrow A_\hbar
\]
fixing the linear generators.
Indeed their star commutator is
$\hbar[-,-]_{\mathfrak g}$.
The induced enveloping-algebra map has associated
graded the identity, and its filtered inverse is $T$.
This algebra map includes the change from the star
product to the PBW product.
Raw PBW symmetrization identifies the modules and
does not by itself identify these products.

The first component of the twisted comparison is
\begin{equation}\label{eq:vol2-finite-formality-component}
 \Phi_{\mathfrak g}(v)
 =T_*\sum_{n\geq0}\frac1{n!}
       \mathcal U_{n+1}
       ((\hbar\pi_{\mathfrak g})^{\otimes n},v).
\end{equation}
It maps $\mathcal P_{\mathfrak g}$ to relative
continuous Hochschild cochains of $A_\hbar$ over
$\mathbb C[[\hbar]]$.
The target differential is the Hochschild differential
of the quantized product.
All terms are defined at each parameter order.
The affine theorem proves the twisted
$L_\infty$ identities on the complexes shifted by $[1]$.

The linear component is a quasi-isomorphism.
Modulo $\hbar$ it is the polynomial HKR map, whose
comparison with the full Hochschild complex is given
by the polynomial diagonal Koszul resolution.
Choose a contracting homotopy of its cone over
$\mathbb C$ and extend it coefficientwise.
The difference between the complete cone differential
and this undeformed one raises $\hbar$ order.
The convergent perturbation series of
Lemma~\ref{lem:vol2-filtered-perturbation}
therefore contracts the complete cone.
This argument also explains why the scaled
Chevalley differential cannot be omitted.

Invariant polynomials have a more explicit comparison.
For $x\in\mathfrak g$, define the formal power series
\begin{equation}\label{eq:vol2-duflo-factor}
 j_{\mathfrak g}^{1/2}(x)
 =
 \det_{\mathfrak g}^{1/2}
 \left(
   \frac{\sinh(\operatorname{ad}_x/2)}
        {\operatorname{ad}_x/2}
 \right),
\end{equation}
with constant term one.
Interpret its Taylor coefficients as
constant-coefficient differential operators on
$\operatorname{Sym}\mathfrak g$.
The Duflo theorem gives the algebra isomorphism
\begin{equation}\label{eq:vol2-duflo-map}
 \operatorname{Duf}_\hbar
 =\operatorname{sym}_\hbar\circ
       j_{\mathfrak g}^{1/2}(\hbar\partial):
 (\operatorname{Sym}\mathfrak g)^{\mathfrak g}[[\hbar]]
 \xrightarrow{\ \sim\ } Z(A_\hbar);
\end{equation}
see \cite{KontsevichDQ}*{Section~8.3.4}.
On a polynomial, only finitely many differential
orders act nontrivially.
Its polynomial Rees formula may consequently be
specialized at $\hbar=1$.
This specialization is not defined on arbitrary
formal power series in $\hbar$.

For $\mathfrak g=\mathfrak h_{\leq d}$, these results
apply to the finite Hamiltonian jets of
Chapter~\ref{ch:vol2-hamiltonian-jets}.
The relative bar--Chevalley comparison, the
$L_\infty$ formality map and the invariant algebra
isomorphism in \eqref{eq:vol2-duflo-map} preserve
different specified operations.
Their conclusions must be used with those operations.

The Hochschild brace action on $(C^\bullet(A_\hbar),A_\hbar)$
can be transported when the required operadic comparison
is supplied. In detail, a homotopy-$E_2$ refinement is
a quasi-isomorphism or a zigzag of quasi-isomorphisms
from the selected complete homotopy-$E_2$ structure
on $\mathcal P_{\mathfrak g}$ to the relative
Hochschild complex, refining the cochain comparison.
A mixed refinement also specifies the open operations
and the coherences between the two colors.

\begin{corollary}[Open--closed structure with a refinement]
\label{cor:vol2-open-closed-refinement}
Given these refinements, the pair
$(\mathcal P_{\mathfrak g},A_\hbar)$ has the transported
homotopy Swiss-cheese action.
Its open product is the enveloping product.
Its closed and mixed operations are those carried by
the specified operadic comparison.
\end{corollary}
\begin{proof}
Transport the two-colored operations through the
supplied equivalence of complete operadic algebras.
The comparison includes the homotopies expressing
substitution and the open--closed identities, so
those identities hold on the transported pair.
\end{proof}

\section{The quadratic and quartic corrections}
\label{sec:vol2-quadratic-quartic-wheels}

For $\mathfrak{sl}_2$, use generators $e,f,h$ with
\[
 [h,e]=2e,\qquad [h,f]=-2f,\qquad [e,f]=h.
\]
In the symmetric algebra set $q=2ef+h^2/2$.
In the Rees enveloping algebra write the generators
as $E,F,H$ and put
$\Omega=EF+FE+H^2/2$.
The relations have an additional factor $\hbar$.
They imply that $\Omega$ is central:
for example,
$[E,EF+FE]=\hbar(EH+HE)$ and
$[E,H^2/2]=-\hbar(EH+HE)$.

\begin{proposition}\label{prop:vol2-duflo-quadratic}
With these normalizations,
\[
 \operatorname{Duf}_\hbar(q)=\Omega+\frac{\hbar^2}{2}.
\]
At $\hbar=1$, its eigenvalue on the irreducible
highest-weight-$n$ module is $(n+1)^2/2$.
\end{proposition}
\begin{proof}
For $x=ue+vf+wh$, the adjoint matrices give
\[
 \operatorname{Tr}(\operatorname{ad}_x^2)=8uv+8w^2.
\]
The first nonconstant term of
\eqref{eq:vol2-duflo-factor} is
$\operatorname{Tr}(\operatorname{ad}_x^2)/48$.
Thus its differential operator on $q$ contributes
\[
 \frac{\hbar^2}{6}
       (\partial_e\partial_f+\partial_h^2)
       (2ef+h^2/2)=\frac{\hbar^2}{2}.
\]
Higher differential orders annihilate $q$, and
symmetrization sends $q$ to $\Omega$.
On a highest-weight vector at $\hbar=1$, $EF$ acts
by $n$, $FE$ by zero and $H^2/2$ by $n^2/2$.
Adding $1/2$ gives $(n+1)^2/2$.
Centrality gives this value on the full irreducible module.
\end{proof}

The next invariant distinguishes a multiplicative
quantization from raw symmetrization.
Use complexified Hermitian generators $J_i$ with
\[
 [J_i,J_j]=i\hbar\epsilon_{ijk}J_k,
 \qquad
 p=x_1^2+x_2^2+x_3^2,\qquad C=\sum_iJ_i^2.
\]
Here $x_i$ are commuting symbols.
The change
$E=J_1+iJ_2$, $F=J_1-iJ_2$, $H=2J_3$
gives $q=2p$ and $\Omega=2C$.
Write $\Delta=\sum_i\partial_{x_i}^2$.
The adjoint eigenvalues are $0,\sqrt p,-\sqrt p$,
so the differential factor becomes
\[
 \frac{\sinh(\hbar\sqrt\Delta/2)}
      {\hbar\sqrt\Delta/2}
 =1+\frac{\hbar^2}{24}\Delta
     +\frac{\hbar^4}{1920}\Delta^2+\cdots.
\]
The square root denotes the displayed power series
in $\Delta$.

\begin{proposition}\label{prop:vol2-duflo-quartic}
In the universal Rees algebra,
\begin{align}
 \operatorname{sym}(p^2)&=C^2-\frac{\hbar^2}{3}C,
       \label{eq:vol2-pbw-quartic}\\
 \operatorname{Duf}_\hbar(p)&=C+\frac{\hbar^2}{4},
       \label{eq:vol2-duflo-compact-quadratic}\\
 \operatorname{Duf}_\hbar(p^2)
       &=\left(C+\frac{\hbar^2}{4}\right)^2.
       \label{eq:vol2-duflo-compact-quartic}
\end{align}
\end{proposition}
\begin{proof}
Antisymmetry of $\epsilon_{ijk}$ proves that $C$ is
central. The commutation relations give
\[
 \sum_j[J_j,[J_j,J_i]]=2\hbar^2J_i,
 \qquad
 \sum_jJ_jJ_iJ_j=(C-\hbar^2)J_i.
\]
The second identity follows by expanding the first.
Symmetrizing the two repeated indices in $p^2$ gives
\[
 \operatorname{sym}(p^2)
 =\frac13\sum_{i,j}
 \bigl(J_i^2J_j^2+J_iJ_jJ_iJ_j+J_iJ_j^2J_i\bigr).
\]
The three sums are $C^2$, $C^2-\hbar^2C$ and $C^2$,
respectively. This proves \eqref{eq:vol2-pbw-quartic}.

Direct differentiation gives
$\Delta p=6$, $\Delta p^2=20p$ and
$\Delta^2p^2=120$.
The quadratic correction is therefore $\hbar^2/4$.
For $p^2$, the full expression is
\[
 C^2-\frac{\hbar^2}{3}C
      +\frac{5\hbar^2}{6}C+\frac{\hbar^4}{16}
 =\left(C+\frac{\hbar^2}{4}\right)^2.
\]
\end{proof}

Thus the product of two operators assigned by $p\mapsto C$
differs from raw symmetrization of $p^2$ by
$\hbar^2C/3$.
This is an operator identity before any state is chosen.
In spin $j$, the eigenvalues of $C$ and
$\operatorname{Duf}_\hbar(p)$ are
$\hbar^2j(j+1)$ and $\hbar^2(j+1/2)^2$.
At $\hbar=1$, evolution by the second operator divided
by two differs from evolution by $C/2$ by $e^{-t/8}$.
Both factors obey the semigroup law, so that law alone
does not choose the vacuum energy.

A multiplicative invariant map with $p\mapsto C$ is
instead
\[
 Q(F(p))=
 \operatorname{Duf}_\hbar
       \bigl(F(p-\hbar^2/4)\bigr).
\]
The shift is an automorphism of the invariant polynomial
algebra. This formula makes the ordering choice explicit;
its relation to an independently prepared state requires
the state comparison as well.

\section{Half-densities and matrix wheels}
\label{sec:vol2-half-density-wheels}

The same formal determinant has a finite-dimensional
geometric meaning. Let $G$ be a local complex Lie group
with unimodular Lie algebra $\mathfrak g$.
For a dimensionless exponential coordinate $X$ near zero,
put $L=\operatorname{ad}_X$.
Left trivialization gives
\[
 d\exp_X=\int_0^1e^{-tL}\,dt
       =e^{-L/2}\frac{\sinh(L/2)}{L/2}.
\]
Normalize the tangent volume and left-invariant volume
to agree at the identity.
Since $\operatorname{Tr}L=0$, the exponential Jacobian
and its normalized square root are
\[
 j(X)=\det\frac{\sinh(L/2)}{L/2},
 \qquad
 \mathfrak j(X)=
 \exp\left(\frac12\operatorname{Tr}
                  \log\frac{\sinh(L/2)}{L/2}\right).
\]
Expanding the scalar series yields
\begin{equation}\label{eq:vol2-log-wheel-coefficients}
 \log\mathfrak j(X)
 =\frac{\operatorname{Tr}L^2}{48}
  -\frac{\operatorname{Tr}L^4}{5760}
  +\frac{\operatorname{Tr}L^6}{362880}
  +O(X^8).
\end{equation}
Each trace is the finite contraction of the adjoint
indices around a wheel.

For $\mathfrak g=\mathfrak{gl}_N$, let
$p_a=\operatorname{tr}(X^a)$, including $p_0=N$.

\begin{proposition}\label{prop:vol2-matrix-wheel-traces}
For every $m\geq0$,
\[
 \operatorname{Tr}_{\mathfrak{gl}_N}(\operatorname{ad}_X)^m
 =\sum_{a=0}^m(-1)^a\binom ma p_{m-a}p_a.
\]
In particular,
\[
 \operatorname{Tr}L^2=2Np_2-2p_1^2,\qquad
 \operatorname{Tr}L^4=2Np_4-8p_1p_3+6p_2^2.
\]
For matrices $X,Y$ of sizes $N,M$, the quadratic part of
$\log\mathfrak j(X\oplus Y)-\log\mathfrak j(X)
 -\log\mathfrak j(Y)$ is
\[
 \frac1{24}
 \left(M\operatorname{tr}X^2+N\operatorname{tr}Y^2
             -2\operatorname{tr}X\operatorname{tr}Y\right).
\]
\end{proposition}
\begin{proof}
On $\operatorname{End}(\mathbb C^N)$, left and right
multiplication commute and their difference is
$\operatorname{ad}_X$.
The matrix-unit basis gives
\[
 \operatorname{Tr}(L_X^aR_X^b)
    =\operatorname{tr}(X^a)\operatorname{tr}(X^b).
\]
Expand $(L_X-R_X)^m$ by the binomial theorem.
The cases $m=2,4$ give the stated polynomials.
Their derivation requires no diagonalizability.
Finally the quadratic logarithm is
$(N\operatorname{tr}X^2-(\operatorname{tr}X)^2)/24$;
subtracting its two block values gives the last formula.
\end{proof}

The off-diagonal matrix blocks are responsible for
this failure of additivity. An adjoint wheel density
is therefore different from a primitive matrix trace.

There is also a convergent product for this
finite-dimensional half-density.
On a sufficiently small neighborhood of zero,
\begin{equation}\label{eq:vol2-half-density-mode-product}
 \mathfrak j(X)=\lim_{K\to\infty}
 \prod_{n=1}^K
       \det\left(1+\frac{L^2}{4\pi^2n^2}\right)^{1/2},
\end{equation}
with branches equal to one at zero.
The logarithms are locally bounded by a constant times
$n^{-2}$, so the product converges locally uniformly.
The Euler product for $\sinh(z/2)/(z/2)$, applied to
the eigenvalues or its matrix power series, identifies
the limit with the determinant above.
If $L$ is real diagonalizable with real eigenvalues,
the partial product $\mathfrak j_K$ satisfies
\[
 0\leq\log\mathfrak j-\log\mathfrak j_K
 \leq\frac{\operatorname{Tr}L^2}{8\pi^2K}.
\]
Indeed $0\leq\log(1+u)\leq u$ for $u\geq0$, and
$\sum_{n>K}n^{-2}\leq1/K$.
For imaginary eigenvalues the scalar factors are
sine factors; the normalized local branch is taken
before their first zeros.
This product has a fixed finite-dimensional adjoint
space. It does not sum traces over infinitely many
independent Lie algebra blocks.

\section{A two-step colour algebra}
\label{sec:vol2-two-step-duflo}

A family in which every wheel vanishes distinguishes
the correction from the underlying comparison.
Let $W$ be finite dimensional and let
$C\in\Lambda^3W^*$.
In dual bases define
\[
 \mathfrak n_C=W\oplus W^*,\qquad
 [x_i,x_j]=\sum_kC_{ijk}y^k,\qquad
 [x_i,y^j]=[y^i,y^j]=0.
\]
The image of every bracket is central, so Jacobi holds
and $\operatorname{ad}_a\operatorname{ad}_b=0$.
The pairing $\kappa(x_i,y^j)=\delta_i^j$, zero on
the other two pairs of summands, is invariant because
$C$ is alternating.

Over $R=\mathbb C[\hbar]$, the Rees enveloping algebra
$A_C$ has the symmetric PBW product
\begin{equation}\label{eq:vol2-colour-pbw-product}
 f*_Cg
 =m\exp\left(
       \frac{\hbar}{2}\sum_{i,j,k}
          C_{ijk}Y^k\,\partial_{X_i}\otimes\partial_{X_j}
       \right)(f\otimes g).
\end{equation}
Every polynomial pair gives a finite sum.
The three bidifferential operators acting on pairs
of slots in a triple tensor product commute, since
their coefficients are central $Y$'s and their
derivatives act only on $X$'s. Both bracketings are
therefore the exponential of the same sum.
The generator commutators are
$[X_i,X_j]_{*_C}=\hbar C_{ijk}Y^k$, as required.

The adjoint coefficient Chevalley complex is
\[
 P_C=R[X,Y]\otimes\Lambda(\xi,\eta),
 \qquad |\xi_i|=|\eta^k|=1,
\]
with differential
\[
 Q\xi_i=QY^k=0,\qquad
 Q\eta^k=-\frac12C_{ijk}\xi_i\xi_j,\qquad
 QX_j=C_{ijk}Y^k\xi_i.
\]
Repeated indices are summed.
These are the Chevalley formulas, and Jacobi gives
$Q^2=0$.
The relative Hochschild cohomology of $A_C$ is
computed by
\[
 Z^p=\operatorname{Hom}_R
       (\Lambda_R^p(\mathfrak n_C\otimes R),A_C).
\]
Its differential uses the commutator action and
$\hbar$ times the bracket faces.
The unnormalized alternating Chevalley-to-bar map
of Chapter~\ref{ch:vol2-hamiltonian-jets}
gives this comparison.

Symmetrization is equivariant:
$[T_a,\operatorname{sym}_\hbar f]
  =\hbar\operatorname{sym}_\hbar(a\cdot f)$.
Consequently
\begin{equation}\label{eq:vol2-colour-beta}
 \beta_p:P_C^p\longrightarrow Z^p,\qquad
 \beta_p=\hbar^p\operatorname{sym}_\hbar
\end{equation}
is a cochain map.
The target differential contributes the extra factor
$\hbar$ needed to pass from degree $p$ to degree $p+1$.

\begin{proposition}\label{prop:vol2-colour-duflo}
Symmetrization identifies the invariant polynomial
algebra with the centre of $A_C$ over $R$.
After extending to $R[\hbar^{-1}]$, $\beta$ induces
an isomorphism of graded cup-product algebras
\[
 H^*(P_C)[\hbar^{-1}]
   \xrightarrow{\ \sim\ }HH_R^*(A_C)[\hbar^{-1}].
\]
The second assertion need not be an integral
isomorphism in positive degree.
\end{proposition}
\begin{proof}
If $f$ is invariant, then
$\sum_{i,k}C_{jik}Y^k\partial_{X_i}f=0$ for every $j$.
The bidifferential operator in
\eqref{eq:vol2-colour-pbw-product} therefore kills
$f\otimes g$ for every $g$, as do all its positive
powers. Hence $f*_Cg=fg$.
Conversely the PBW commutator of a linear generator
with $f$ is $\hbar$ times its adjoint derivative.
The PBW module has no $\hbar$-torsion, so centrality
is equivalent to invariance. This proves the
integral degree-zero assertion.

After inverting $\hbar$, identify $A_C$ with
$U(\mathfrak n_C)\otimes R[\hbar^{-1}]$ by
$T_a\mapsto\hbar a$.
Evaluate a degree-$p$ cochain on the rescaled
arguments $\hbar^{-1}T_a$; this contributes
$\hbar^{-p}$ and cancels the factor in $\beta_p$.
The target becomes the ordinary adjoint Chevalley
complex. Under these identifications, $\beta$ is
symmetrization composed with the source automorphism
$f(X,Y)\mapsto f(\hbar X,\hbar Y)$.
This automorphism commutes with $Q$ and preserves
the cup product.
Every adjoint product of length at least two is
zero, so the determinant
\eqref{eq:vol2-duflo-factor} equals one.
The cohomological Duflo theorem of Pevzner and
Torossian \cite{PevznerTorossianDuflo}*{Theorem~2.5}
then proves the cup-product assertion in every
degree. The polynomial argument extends over the
Laurent coefficient ring.

For the last assertion take $C=0$.
Both cochain differentials are zero, and $\beta_p$
is multiplication by $\hbar^p$ on a nonzero free
module whenever $0<p\leq\dim\mathfrak n_C$.
Its cokernel has $\hbar$-torsion and its
specialization at $\hbar=0$ is zero.
\end{proof}

There are algebraic product homotopies after this
localization, although the theorem does not
choose a field-theoretic one.
Here is an explicit way to specify the first.
Let $X=P_C\otimes P_C$ and $Y=Z$, now over
$R[\hbar^{-1}]$, and write
\[
 D=\beta\,m_{P_C}-m_Z(\beta\otimes\beta):X\longrightarrow Y.
\]
The cup product on $Z$ uses the usual Chevalley
shuffle diagonal and the coefficient product.
Choose contractions $(i_X,p_X,h_X)$ and
$(i_Y,p_Y,h_Y)$ onto cohomology.
They exist here: the preceding rescalings identify
the complexes with extensions of complexes of
complex vector spaces, which split into cohomology
and contractible pairs.
Multiplicativity on cohomology says $p_YDi_X=0$.
Then
\begin{equation}\label{eq:vol2-colour-cup-homotopy}
 H=h_YD+i_Yp_YDh_X
 \quad\hbox{satisfies}\quad d_YH+Hd_X=D.
\end{equation}
To verify it, use that $D$ is a cochain map and
expand the two contraction identities.
The first summand gives $(1-i_Yp_Y)D$.
The second gives $i_Yp_YD$, since
$p_YDd_X=0$ and $p_YDi_X=0$.
Their sum is $D$.
This construction depends on the splittings.
Compatibility with braces, geometric collisions
and a prescribed two-point operation is further
structure.

\section{A mode completion that excludes the finite normalization}
\label{sec:vol2-wheel-mode-obstruction}

Finite Duflo maps do not automatically assemble along
projections that discard Lie algebra blocks.
For each $i\geq1$, let $\mathfrak g_i$ be a copy
of $\mathfrak{sl}_2$ with generators $E_i,F_i,H_i$.
Different blocks commute. Put
$\mathfrak g_{\leq N}=\bigoplus_{i=1}^N\mathfrak g_i$
and define
\[
 A_{\mathrm{mode}}=
 \varprojlim_{N,M}
       \operatorname{Sym}(\mathfrak g_{\leq N})
                  \otimes\mathbb C[\hbar]/(\hbar^M),
\]
\[
 U_{\mathrm{mode}}=
 \varprojlim_{N,M}
       U_\hbar(\mathfrak g_{\leq N})/(\hbar^M).
\]
The transition maps kill the omitted blocks.
Each finite quotient has the discrete linear topology,
and the limits have their inverse-limit topologies.

\begin{proposition}\label{prop:vol2-mode-duflo-obstruction}
There is no continuous linear map
$A_{\mathrm{mode}}\to U_{\mathrm{mode}}$ whose
restriction to invariant polynomials in every finite
set of blocks is the Duflo map with the normalization
of \eqref{eq:vol2-duflo-map}.
\end{proposition}
\begin{proof}
The invariant
$c_i=H_i^2/2+2E_iF_i$ tends to zero in
$A_{\mathrm{mode}}$: every fixed block projection
eventually kills it.
Proposition~\ref{prop:vol2-duflo-quadratic} prescribes
its image as
\[
 \operatorname{sym}_\hbar(c_i)+\frac{\hbar^2}{2}.
\]
Project to the first block and reduce modulo
$\hbar^3$. For every $i>1$ the PBW term is zero,
while $\hbar^2/2$ is nonzero.
The proposed images do not tend to zero, contradicting
continuity.
\end{proof}

The obstruction fixes both the topology and the
normalization on every finite block.
It does not exclude a different coefficient domain
or a comparison with a specified renormalization.
For example, using a parameter $q_i$ of increasing
positive weight as the $i$th Rees parameter replaces
the scalar correction by $q_i^2/2$.
Only finitely many such corrections survive at any
fixed weight. The next construction makes that
coefficient domain precise.

\section{Positive weights and mixed coefficients}
\label{sec:vol2-weighted-mixed-formality}

Let $A=\mathbb C[x_1,x_2,\ldots]$, where each polynomial
uses finitely many variables.
Write $T_{\mathrm{fin}}^{p-1}$ for finite sums of
polynomial $p$-vectors and
$D_{\mathrm{fin}}^{p-1}$ for finite sums of operators
\[
 (f_1,\ldots,f_p)\longmapsto
 a(x)\,\partial^{\alpha_1}f_1\cdots\partial^{\alpha_p}f_p.
\]
Every multi-index has finite support.
In degree $-1$, both spaces are $A$.
These are the shifted deformation complexes:
a $p$-vector or $p$-ary cochain has degree $p-1$.
Use the Gerstenhaber insertion bracket on $D_{\mathrm{fin}}$.
If $m$ is ordinary multiplication, its differential is
$b_0=[m,-]_G$.
For the Schouten bracket use
\[
 [X_1\wedge\cdots\wedge X_p,f]_{\mathrm{Sch}}
 =\sum_i(-1)^{i-1}X_i(f)
          X_1\wedge\cdots\widehat X_i\cdots\wedge X_p.
\]

Take a mixed coordinate domain $U=I^a\times D^b$,
with $I$ an open real interval and $D$ a complex disk.
Let $K(U)$ be the algebra of smooth real de Rham
forms and complex $(0,*)$-forms on this product,
with
$d_K=d_{\mathbb R^a}+\bar\partial$
and the total-complex signs.
Its coefficients are arbitrary smooth functions
on the product, not just finite sums of separated
functions. Give it the compact-open $C^\infty$ topology.

Give each finite polynomial operator space the fine
bornology: a bounded set is contained and bounded
in a finite-dimensional subspace.
In $K(U)\widehat\otimes T_{\mathrm{fin}}$ and
$K(U)\widehat\otimes D_{\mathrm{fin}}$, this means
the direct sum of copies of $K(U)$ indexed by a
polynomial operator basis; a bounded set uses only
finitely many basis directions.
Field derivatives act on the $x_i$ only.
No spacetime differential operator or product of
distributions is included in this definition.

Choose degree-zero central parameters $q_i$ with
positive integer weights $w_i$, only finitely many
of which lie below any fixed bound.
Set
\[
 \operatorname{wt}(q^\alpha)=\sum_iw_i\alpha_i,
 \qquad
 \mathscr R_{<M}
   =\mathbb C[q_1,q_2,\ldots]/
       (q^\alpha:\operatorname{wt}(q^\alpha)\geq M),
 \qquad
 \mathscr R=\varprojlim_M\mathscr R_{<M}.
\]
Each $\mathscr R_{<M}$ is finite dimensional.
For the parameter-completed tensor define
$V\widehat\otimes\mathscr R
 =\varprojlim_M(V\otimes\mathscr R_{<M})$ and put
\[
 \widehat T(U)=
   (K(U)\widehat\otimes T_{\mathrm{fin}})
                 \widehat\otimes\mathscr R,\qquad
 \widehat D(U)=
   (K(U)\widehat\otimes D_{\mathrm{fin}})
                 \widehat\otimes\mathscr R.
\]
The filtration $F^M$ is the kernel of projection
to parameter precision $M$.
Every coefficient has finite field support,
although that support can increase with weight.

Finite polynomial differential graphs act on these
spaces. On a bounded input set, only finitely many
coefficient and derivative directions occur.
A graph consequently has finitely many index
assignments. Its spacetime coefficients are products,
and Leibniz's rule bounds their derivatives on each
compact subset by finitely many products of input
seminorms. Thus the graph map is bounded.
A series whose $n$th term raises weight by at least
$n$ has only finitely many terms at precision $M$,
provided that each term has finitely many graphs.
Its compatible finite-precision maps define a
continuous map on the completion.

Extend the affine Taylor maps over the mixed
coefficients by
\begin{equation}\label{eq:vol2-mixed-formality-sign}
 \mathcal U_n^K(k_1\gamma_1,\ldots,k_n\gamma_n)
 =
 (-1)^{\sum_{i<j}|\gamma_i||k_j|
              +(1-n)\sum_i|k_i|}
 k_1\cdots k_n\,\mathcal U_n(\gamma_1,\ldots,\gamma_n).
\end{equation}
The first exponent moves the coefficients to the
left. The second is the tensor sign for a map of
degree $1-n$.
After suspension, the Taylor maps have degree zero.
Their coefficient extensions commute with the
total coderivations because $d_K$ is a derivation.
Undoing suspension proves all the extended
$L_\infty$ identities with the displayed sign.
The finite graph formulas are stable under
adjoining unused field coordinates.

Let $\Pi\in F^1\widehat T(U)$ be a bivector of
mixed-form degree zero satisfying
$d_K\Pi=0$ and $[\Pi,\Pi]_{\mathrm{Sch}}=0$.
Independent finite Poisson blocks $\pi_i$ give the
example $\Pi=\sum_iq_i\pi_i$.

\begin{theorem}\label{thm:vol2-weighted-relative-formality}
The formulas
\begin{align}
 B_\Pi&=\sum_{n\geq1}\frac1{n!}
              \mathcal U_n^K(\Pi^{\otimes n}),
       &m_\star&=m+B_\Pi,
       \label{eq:vol2-weighted-star}\\
 \mathcal F_r(v_1,\ldots,v_r)
       &=\sum_{n\geq0}\frac1{n!}
           \mathcal U_{n+r}^K
                (\Pi^{\otimes n},v_1,\ldots,v_r),
       &&r\geq1
       \label{eq:vol2-weighted-twist}
\end{align}
define an associative product commuting with $d_K$
and a filtered continuous $L_\infty$ quasi-isomorphism
\[
 (\widehat T(U),d_K+[\Pi,-]_{\mathrm{Sch}})
 \longrightarrow
 (\widehat D(U),d_K+[m_\star,-]_G).
\]
The maps commute with restriction in $U$.
The assertion is about the unreduced mixed complexes.
\end{theorem}
\begin{proof}
At precision $M$, the twisting sums have $n<M$.
The input coefficients and the relevant coefficients
of $\Pi$ use a common finite set of field variables.
The calculation is therefore the finite affine
formality calculation with coefficients in $K(U)$
and signs \eqref{eq:vol2-mixed-formality-sign}.
The Maurer--Cartan equation gives
\[
 (d_K+[m,-]_G)B_\Pi+\frac12[B_\Pi,B_\Pi]_G=0.
\]
Since $B_\Pi$ is binary and has mixed-form degree
zero, this equation gives separately
$d_KB_\Pi=0$ and $[m_\star,m_\star]_G=0$.

For the higher maps, conjugate the source symmetric
coalgebra codifferential by multiplication with
$\exp(\Pi)$ and the target one by $\exp(B_\Pi)$.
The coefficient with $r$ free inputs is
\eqref{eq:vol2-weighted-twist}.
All coefficients and identities are finite at
the specified precision. Compatibility in $M$
proves them on the inverse limit.

The untwisted finite-direction HKR map is a
quasi-isomorphism by the filtered colimit of the
finite polynomial HKR comparisons.
Its cone is an acyclic complex of complex vector
spaces. Choose a square-zero contracting
homotopy $H_0$ by splitting it into contractible
pairs. Every linear map between fine bornological
spaces is bounded.
Extend $H_0$ over $K(U)$ with its degree-minus-one
tensor sign. The two terms involving $d_K$ cancel,
so this contracts the unreduced mixed cone.

Write the cone differential of $\mathcal F_1$ as
$d_0+\delta$. Every term of $\delta$ raises
parameter weight.
The homotopy
\[
 H=H_0(1+\delta H_0)^{-1}
   =H_0\sum_{j\geq0}(-\delta H_0)^j
\]
is a finite sum at every precision.
The perturbation identity gives
$(d_0+\delta)H+H(d_0+\delta)=1$.
This proves the quasi-isomorphism and continuity.
Finally the graph maps use only field derivatives
and pointwise products of mixed coefficients,
which commute with restriction in $U$.
\end{proof}

Both sign contributions in
\eqref{eq:vol2-mixed-formality-sign} can occur.
On an interval with coordinate $t$, take
$\pi=\partial_u\wedge\partial_v$ and $\Pi=h\pi$,
where $h$ has weight one.
The constant-bivector normalization gives
$\mathcal U_2(\pi,\pi)(u^2,v^2)=1$.
The order-$h$ term of $\mathcal F_1(t\pi)$ on
these two functions is $ht$.
Its mixed differential is $h\,dt$.
In $\mathcal U_2^K(\pi,dt\,\pi)$, the input-crossing
sign and the degree-minus-one map sign cancel,
giving precisely this value.

The target remains the completed finite-direction
polydifferential complex.
For example,
$\sum_i x_i\partial_{x_i}$ acts on every polynomial
and is bounded for the fine bornology.
It is not a finite-direction derivation.
It is not a Hochschild boundary either, since
the nullary Hochschild differential of a
commutative algebra is zero.
Thus a comparison with the full bounded
Hochschild complex requires a larger source.

\section{A variable Poisson coefficient}
\label{sec:vol2-variable-poisson-test}

A small calculation shows why both the target differential
and the derivative corrections matter.
On $\mathbb C[x,y]$, take
$\pi=x\partial_x\wedge\partial_y$ and use one
parameter $h$.
The first coefficient of the star product is
\[
 B_1(f,g)=\frac{x}{2}(f_xg_y-f_yg_x).
\]
The Schouten and Hochschild formulas give
\[
 [\pi,x]_{\mathrm{Sch}}=x\partial_y,\qquad
 [\pi,\partial_x]_{\mathrm{Sch}}
       =-\partial_x\wedge\partial_y,
\]
\[
 [m_\star,x]_G(y)=hx+O(h^2),\qquad
 [m_\star,\partial_x]_G(x,y)=-h/2+O(h^2).
\]
For the last expression, expand
$(\partial_xx)\star y+x\star(\partial_xy)
 -\partial_x(x\star y)$.
The first-order term is $-h/2$.

There cannot be a cochain map reducing to HKR from
$(T_{\mathrm{fin}}[[h]],h[\pi,-]_{\mathrm{Sch}})$
to the undeformed target
$(D_{\mathrm{fin}}[[h]],b_0)$.
Such a map sends $x$ to a nullary cochain, on which
$b_0$ is identically zero.
Its order-$h$ chain equation would be
$0=h\,\operatorname{HKR}([\pi,x])=hx\partial_y$.
Higher coefficients cannot change that equation.

For the second product term, choose the affine gauge
with
\begin{equation}\label{eq:vol2-variable-poisson-second-term}
 \begin{split}
 B_2(f,g)={}&\frac{x^2}{8}
       (f_{xx}g_{yy}-2f_{xy}g_{xy}+f_{yy}g_{xx})\\
 &+\frac{x}{12}
       (f_xg_{yy}+f_{yy}g_x-f_{xy}g_y-f_yg_{xy}).
 \end{split}
\end{equation}
This is the second-order affine formula
\cite{KontsevichDQ}*{Section~1.4.2}, normalized so
that the first commutator is $h\{-,-\}$.

\begin{proposition}\label{prop:vol2-variable-poisson-associator}
The product $f\star g=fg+hB_1(f,g)+h^2B_2(f,g)$
is associative modulo $h^3$.
In particular,
\[
 (x\star y)\star y=x\star(y\star y)
   =xy^2+hxy+\frac{h^2x}{6}\pmod{h^3}.
\]
Omitting the $x/12$ term in
\eqref{eq:vol2-variable-poisson-second-term}
leaves associator $h^2x/4$ on $(x,y,y)$.
\end{proposition}
\begin{proof}
Leibniz's rule cancels the first-order associator.
The second coefficient on $(f,g,k)$ is
\[
 \begin{split}
 &B_2(f,g)k+B_1(B_1(f,g),k)+B_2(fg,k)\\
 &\hspace{1cm}
   -fB_2(g,k)-B_1(f,B_1(g,k))-B_2(f,gk).
 \end{split}
\]
Using only the $x^2/8$ term and expanding each
derivative by Leibniz gives
$x(f_xk_y-f_yk_x)g_y/4$.
The four terms with coefficient $x/12$ give its
negative. This proves the assertion for arbitrary
polynomials.
For the displayed triple,
$B_2(xy,y)=-x/12$ and $B_2(x,y^2)=x/6$.
The iterated $B_1$ term adds $x/4$ on the left,
giving $x/6$ on each side.
Without the derivative correction, the left side
has $x/4$ and the right side has zero in order two.
\end{proof}

The parameter also affects Lie compatibility.
After inverting $h$, multiplication by $h^{p-1}$
on $p$-vectors identifies the differentials
$[\pi,-]$ and $h[\pi,-]$ and preserves the
Schouten bracket: on a bracket of a $p$-vector and
a $q$-vector both powers are $h^{p+q-2}$.
It sends the function $1$ to $h^{-1}$.
The unit-preserving scaling $h^p$ instead preserves
the bracket if the source bracket is multiplied
by $h$. The two rescalings answer different
comparison questions.

\section{Comparing quantizations and continuous operations}
\label{sec:vol2-continuous-wheel-comparison}

Two quantizations on the same coefficient module
must be compared in their allowed operator space.
Suppose associative products $\mu=m+O(h)$ and
$\mu'=m+O(h)$ agree modulo $h^N$, for $N\geq1$.
Assume their coefficient domains and differentials
have been identified, and put
$\Delta_N=h^{-N}(\mu'-\mu)\bmod h$.
Subtracting associativity in order $h^N$ gives
$b_0\Delta_N=0$.
A permitted degree-zero unary operator $D$ gives
\[
 \mu'=(1+h^ND)^{-1}
        \mu\bigl((1+h^ND)\otimes(1+h^ND)\bigr)
                      \pmod{h^{N+1}}
\]
if and only if
\begin{equation}\label{eq:vol2-unary-gauge-equation}
 \Delta_N=b_0D.
\end{equation}
Indeed expansion of the right side gives
\[
 \mu+h^N\bigl(m(D\otimes1)+m(1\otimes D)-Dm\bigr).
\]
If the mixed differential $d_K$ is fixed, the
same comparison also requires $[d_K,D]=0$.
A primitive in a larger space of operators does
not solve the equation in the prescribed domain.
For a commutative product $b_0D$ is symmetric in
its inputs, so a nonzero antisymmetric first-order
difference cannot be removed this way.
For example,
$b_0(\partial_x^2)(f,g)=-2f_xg_x$.

The Hamiltonian jet system has a different
completion from the finite-direction weighted
model. The two-index construction of
Chapter~\ref{ch:vol2-continuous-cochains}
gives $\mathcal P_{\mathrm{uc}}$ and its cochain
comparison with $C^\bullet_{\mathrm{uc}}(A/R,A)$.
A geometric polyvector realization must first be
compared with that complex in its stated topology.
The compatible-jet criterion of
Theorem~\ref{thm:vol2-compatible-jets} records the
additional transitions and limit identifications
needed for a proposed $\mathcal P_\infty$.

The further question is a continuous
homotopy-$E_2$ quasi-isomorphism into this
two-cutoff Hochschild recipient, with an invariant
Duflo map and a specified renormalization of the
wheels. A negative answer would require a brace
obstruction that persists under every allowed
renormalization. The fixed-normalization
counterexample of
Proposition~\ref{prop:vol2-mode-duflo-obstruction}
does not have that quantifier.

At the binary product level a proposed cochain
map $\Phi$ already needs a homotopy satisfying
\[
 dH(u,v)+H(du,v)+(-1)^{|u|}H(u,dv)
   =\Phi(uv)-\Phi(u)\smile\Phi(v).
\]
The brace substitutions require compatible higher
homotopies on the same completed coefficient
domain. A finite invariant algebra isomorphism
does not specify them.
The mixed-chain traces, cyclic orientation and
finite wheel operators constructed above provide
separate parts of such a comparison.
To identify their action with an independently
defined field theory, one must also preserve
insertion supports, singular coefficients,
collision substitutions and the chosen state.

